\section{Introduction}\label{sec:introduction125}

An investment decision depends on the continuation values of the actions
that remain feasible, on the risk of their future consequences, and on the
cost of obtaining those values. A small average prediction error need not
identify a good investment. An accurate certificate for a fixed policy does
not by itself improve that policy. Neural Bellman Operators address this
economic decision problem by constructing a continuation critic together
with a feasible actor for a specified recursive objective. Their numerical
content must be assessed at the level of that complete decision service.

Our application is a constrained investment economy with a nonlinear
transition, a state-dependent capacity, a continuous innovation, and a nested
tail-risk preference. The primitive Bellman objective, rather than a fitted
surrogate or a discretized alternative economy, defines optimality. We ask
three related questions. What information about the next action is already
contained in a paid Bellman calculation? When should a service obtain a new
action observation or a more accurate evaluation of an existing one? Can a
neural continuation provide enough economic value to pay for its acquisition,
training, qualification and deployment relative to strong conventional
methods with the same information?

The first contribution is an exact finite description of the original-action
information in a complete scalar Bellman support record. Existing
qualification calculations return a global quadratic minorant as well as
an interval for the value at the queried action. Compressing this record to
a value and a derivative interval discards restrictions on other actions.
We retain every quadratic minorant. After removing the known strong
convexity, each minorant becomes an affine obstacle. Adjacent smooth-convex
interpolation conditions, together with a bounded-domain conjugate
inequality at each node and obstacle, are necessary and sufficient for
compatibility with all the original records. The proof constructs a
continuous piecewise-affine derivative realizing the information. A finite
convex program therefore gives the exact worst compatible action contrast
under this explicitly stated information interface. A rational residual
correction validates a bound even when an optimization routine supplies an
inexact candidate. This is a completion of the information used by the
Bellman certificate, with no claim that convex interpolation or dual
residual certification is itself new.

The second contribution links action localization to the complete cost of
continuous-risk information. The inherited midpoint rule already bounds
bracket contraction. We make explicit the support precision required for a
declared contraction, the sum of all refinement costs, the constrained
action-rounding allowance, and the separate terminal requirements for
original-policy regret and absolute objective width. Under a stated
tunable-oracle cost assumption, the resulting precision schedule has a
finite work bound; the assertion does not presume that a fixed experimental
precision ladder reaches the required level. When the available information
cannot certify a change, the algorithm preserves the last qualified feasible
policy and charges the work that was attempted.

These developments retain the original neural construction. Its critic is
built backwards from its own continuation and its actor carries a feasible
action witness. The operator foundations establish finite-horizon value
existence and uniqueness, measurable decisions, stability and implementation
allowances. One-sided and centered-risk accounts connect local errors to the
original recursive objective. The supplement supplies complete proofs and
the conditional bridge to stochastic control and nonlinear recursive
utility. The Brownian bridge requires its own well-posed backward equation,
regularity and stopping-boundary conditions. A result under expectation,
nested AVaR, or a value-dependent recursive driver retains that particular
objective and information filtration.

The numerical study asks whether the new acquisition service improves the
decision after paying for a strong original-model control. It uses all
56 previously acquired actors, covering neural ReLU, exactly matched
ordinary hinge, quadratic, spline, polynomial and grid representations,
and includes direct online and uniform-action dynamic-programming controls.
All economic comparisons retain the same continuous innovation law,
original future controller and service tolerances. Four states and the
complete method catalogue were fixed in a remotely recorded source
boundary before the first new evaluation. This deterministic study measures
action information and complete work on that specified census. Earlier
independent random-state evidence and its simultaneous inference remain
separately identified. Repeated process clocks are descriptive measurements,
not additional economic observations.

The strongest historical comparisons constrain the interpretation of the
new results. Exact ReLU and matched hinge encodings return the same
decisions. Previous complete-cost neural adoption domains were empty, and
traditional quadratic methods often won. We preserve these results and
their failed or unresolved computations. The present support theorem
applies equally to conventional and neural proposals; its rational strict
witness also has a cheap direct-support comparator. It demonstrates that
discarded information can affect a certified economic gain, not a unique
neural advantage. The original-economy experiment implements the separately
specified acquisition service; it is not an executed performance test of
the new complete-support convex program.

The contribution sought by Neural Bellman Operators remains a competitive
method for recursive economic decisions. The current results identify
precisely which parts of that contribution are established: feasible
own-future construction, original-objective certificates, exact action
information, and an accountable mechanism for acquiring further precision.
Whether a learned neural representation can deliver a positive complete
economic advantage over equally informed conventional representations
remains an empirical and theoretical question, addressed by the comparisons
and explicit adoption conditions rather than assumed from the architecture.

The article first specifies the operator and original economy. It then
develops the paid-support and complete-information results, the precision
and work account, and the current economic experiment. The technical
supplement contains the complete proof chains. A separate development
companion preserves the earlier models, statements, executions, negative
results and corrections, with source identities and a content-retention map.
